\section{The fragmentation problem}
\label{sec:introduction}

A difficult result in mathematical physics is rarely proved in the variables
in which it receives its final physical interpretation.  Gauge potentials
make local dynamics tractable but contain redundancy.  Gauge-invariant
observables represent physical quantities more directly but can obscure local
coordinates.  A lattice regulator replaces an infinite-dimensional field
problem by finite group variables while breaking continuous spacetime
structure.  Renormalization reorganizes those variables by scale.  Euclidean
correlation functions expose positivity and probabilistic control.  A
Hamiltonian representation finally turns long-time behavior into a spectral
statement.  Each description is useful because it displays some structure
and suppresses other structure.

This plurality creates a logical problem that is sharper than the familiar
observation that models are partial.  Suppose an estimate is proved after
gauge fixing, coarse graining, finite-volume restriction, or passage to a
Euclidean measure.  Which part of the estimate belongs to the physical theory,
and which part belongs to the chosen representation?  What must be shown
before a local contraction becomes a statement about continuum
correlations, or before exponential Euclidean decay becomes a Hamiltonian
mass gap?  A general appeal to equivalence is inadequate because the relevant
maps are often irreversible, approximate, indexed by singular limits, or
equivalent only on a restricted algebra of observables.

The philosophy of physics already provides indispensable pieces of the
answer.  Structural realism asks which structure merits realist commitment
through theory change \cite{Worrall1989,VanFraassen2006,Roberts2011}.
Work on gauge symmetry distinguishes reduction, sophistication, and surplus
structure, including the role that unreduced variables can play in
representing local fields \cite{Belot2003,Dewar2019,NguyenTehWells2020}.
Accounts of theoretical equivalence and duality separate formal isomorphism
from interpretation \cite{Coffey2014,Lupher2018,DeHaroButterfield2018}.
Functorial quantum field theory makes covariance and inter-model maps
mathematically precise \cite{BrunettiFredenhagenVerch2003,Fewster2015}.
Renormalization-group and effective-field-theory scholarship shows that scale
transformation has several technically and epistemically different forms
\cite{WilsonKogut1974,Butterfield2011,Rivat2021,KoberinskiFraser2023}.
Recent work on perspicuous representation further argues that physical
formulations are not selected by convention alone \cite{North2026}.

These approaches mostly ask what theories, models, structures, or
perspectives are, when descriptions are equivalent, or why one formulation
is more revealing.  The present question is narrower and proof-directed:

\begin{quote}
When does a property proved in one representation survive a chain of
nonexact transformations and count as established in a limiting physical
representation?
\end{quote}

We answer by associating each transformation with a \emph{certificate
profile}.  A profile is not a global label such as ``faithful.''  It records
which property channels are exact, which have observable-dependent error
moduli, which are unavailable in the target variables, which are recovered by
a separate theorem, and over which cutoff, volume, scale, distance, and
observable quantifiers the certificate is uniform.  We call the resulting
analysis \emph{proof-bearing representational dynamics}.  The qualifier is
essential: ``representational dynamics'' already names temporal changes of
internal representations in cognitive science and machine learning.  Here it
means transformations among scientific descriptions together with proof
objects for claim transport.

The paper makes four contributions.  First, it defines a typed representation
package containing a redundancy groupoid, a physical observable algebra,
positive states, dynamics or correlations, scale/limit data, and certified
properties.  Its morphisms distinguish exact preservation, approximation,
forgetting, and recovery.  Their error moduli compose by typed pullback rather
than by collapsing every property into one number.

Second, we prove a recovery-to-spectrum theorem.  Consider positive
semigroups generated by $H_i\geq0$ and a target $H\geq0$.  If recovered
quadratic forms converge on a centered core that is dense in the target
vacuum complement, and if every regulated form obeys one exponential rate
$m>0$ uniformly in the index, then
\[
  \spec(H)\cap(0,m)=\varnothing.
\]
The mechanism is positivity of the spectral measure, followed by density of
the core.  The theorem is elementary as spectral analysis; its role is to
make four prerequisites of a physical inference inseparable: positivity,
correlation recovery, observable completeness, and uniform physical scale.
It therefore exceeds a Boolean certificate conjunction or a scalar
error-summing rule.

Third, three exact countermodels establish the independence of those
prerequisites.  A nonseparating core can miss a lower-energy sector.  A
sequence of finite systems can have a positive gap at every index while the
gaps converge to zero.  A recovery defect can vanish at each fixed volume but
remain one after the infinite-volume limit.  These are not pathological
decorations; they reproduce the logical forms of common shortcuts in
constructive field theory.

Fourth, the Yang--Mills existence and mass-gap problem provides a sustained
case rather than a motivational slogan.  We track locality, gauge invariance,
positivity, scale dependence, observable completeness, and spectrum from
lattice link variables through gauge quotient and block transformation to
continuum Schwinger data, Osterwalder--Schrader (OS) reconstruction, and a
Hamiltonian.  Four concrete failures of proposed lattice estimates are
classified as losses of quantifier domain, boundary state, adaptive history,
and normalization.  The classification is methodologically effective because
each loss selects a different mathematical repair.

The official problem asks for a nontrivial quantum Yang--Mills theory on
$\mathbb R^4$ and a positive Hamiltonian mass gap
\cite{JaffeWitten2000}.  Nothing here constructs that continuum theory or
proves the uniform infrared estimate needed to apply the central theorem.
The Yang--Mills contribution is a precise interface theorem and a conservative
frontier diagnosis.  The formal result is a conditional representation-
transfer theorem; the physical premises remain open.

Section~\ref{sec:levels} separates ontology, semantics, epistemology, and
proof methodology.  Sections~\ref{sec:calculus}--\ref{sec:countermodels}
develop the calculus, theorem, and countermodels.  Sections~\ref{sec:ymchain}
and~\ref{sec:failures} analyze Yang--Mills.  Sections~\ref{sec:ontology}
and~\ref{sec:related} locate the framework philosophically, after which
Sections~\ref{sec:objections} and~\ref{sec:discussion} address objections and
research consequences.  Appendices contain complete definitions, proofs,
the property-transport table, and exact finite reproducibility artifacts.
